%% spread.tex -- a two-sided document: the garden follows the outer margin from page
%% to page, and on the last page it is grown from another theme. LuaLaTeX, one pass:
%%   lualatex spread.tex
\documentclass[11pt,a4paper,twoside]{article}
\renewcommand\rmdefault{cmr}\renewcommand\sfdefault{cmss}\renewcommand\ttdefault{cmtt}
\usepackage[T1]{fontenc}
\usepackage{amsmath,amsthm}
\usepackage[a4paper,twoside,inner=22mm,outer=68mm,top=32mm,bottom=30mm,headsep=8mm]{geometry}
\usepackage{fancyhdr}
\usepackage[grown,sections]{florilegium}

\pagestyle{fancy}\fancyhf{}
\fancyhead[RO,LE]{\small\textsc{Two sides, two gardens}}
\fancyfoot[RO,LE]{\small\thepage}
\renewcommand\headrulewidth{0pt}
\newcommand\code[1]{\texttt{\small #1}}
\theoremstyle{definition}
\newtheorem*{remark}{Remark}

\begin{document}
\section{Both sides of the page}

This document is printed on both sides of the paper. Its outer margin is the wide one, so
on this page, a right-hand page, the margin is on the right; turn the page and it is on the
left. The garden follows it without being told: it grows in whichever side margin is the
wider, and in a two-sided document that changes with every page. The document asks for
nothing more than the usual two-sided layout:
\begin{center}\code{\textbackslash documentclass[twoside]\{article\}}\\
\code{\textbackslash usepackage[twoside,inner=22mm,outer=68mm]\{geometry\}}\end{center}

\subsection{Where the margin is}

In a two-sided document \LaTeX\ moves the text block by \code{\textbackslash oddsidemargin}
on odd pages and by \code{\textbackslash evensidemargin} on even ones. When a page is
shipped out, the package reads the same one, finds the left and right edges of the text,
and compares the two margins. If the left one is the wider, nothing changes: this is the
page the garden has always grown on. A one-sided document with a wide right margin is
treated like a right-hand page, on every page.

\subsection{A garden in a mirror}

If the right margin is the wider, the page is laid out as if it were seen in a mirror. A
point at $x$ centimeters from the left edge of the paper moves to $w - x$, where $w$ is the
paper's width, so the text block's two edges trade places and the wide margin is on the
left again. The growth code is not told any of this. It finds the column, the footer band
and the free space under the text where it always does, grows its stems from a root in the
bottom-left corner, and fills the holes one by one.

Only the last step changes. Every point of every path is mapped from layout centimeters to
PDF points by one affine transform before it is written, and on a mirrored page that
transform reflects:
\[
  \begin{pmatrix} x' \\ y' \end{pmatrix}
  = \frac{72}{2.54}\begin{pmatrix} -1 & 0 \\ 0 & 1 \end{pmatrix}
    \begin{pmatrix} x \\ y \end{pmatrix}
  + \frac{72}{2.54}\begin{pmatrix} w \\ 0 \end{pmatrix}.
\]
Everything goes through it: each petal and vein, the snail, which now crawls the other way,
the bird, which faces the other way, and the copper ivy leaf in the header, whose stem
runs out into the head rule toward the inside of the page. Line widths are given in
absolute points and are not affected, and neither are the fills: a reflection reverses
the direction of every outline, which turns each winding number into its negative, and
the nonzero rule only asks whether it is zero.

PDF could do the reflection itself: one \code{cm} operator in front of the drawing would
reflect everything after it. The garden does not use \code{cm} anywhere. Each point is
transformed in Lua before it is written, which is what keeps line widths in absolute
points while a motif is scaled, and what lets a primitive called inside a TikZ scope read
pgf's matrix and draw under it. The reflection is one more matrix in that chain, and the
page's content stream gains no operator it did not have before: only its numbers differ.

\begin{remark}
The ivy leaf is the one place where this went wrong. It reports where its stem ends, so
that the head rule can continue from that point, and it reports it in page coordinates,
which are already reflected. The rule is then drawn in layout coordinates and reflected a
second time, so on the first mirrored pages it started in the wrong margin. The fix is
one line: reflect the point back before drawing the rule.
\end{remark}

\subsection{Same seed, same garden}

A page's garden still comes from its seed, and the mirror changes none of the random
draws. A mirrored page is therefore exactly the reflection of the page that the same seed
would give in a left margin, petal for petal. On the last page, where the text ends early,
the space under it is mirrored as well, so the branches that arch into it grow out of the
garden on the correct side.

The flowers beside the section headings work as before. A heading leaves an invisible mark
in the text, the page box is searched for it at shipout, and the mark gives only a height;
the side does not matter, so the copper flower lands beside the heading in either margin.

The footer runs the other way too. Its stem leaves the root in the outer corner and
runs along the bottom of the page toward the binding, thinning to a copper hairline under
the text, and the header's rule runs from the ivy leaf in the outer margin across the top
of the text block. The header's words and the page number are set by \code{fancyhdr}, on
the outside of each page, with \code{\textbackslash fancyhead[RO,LE]}.

\begin{remark}[Two pages, one picture]
Seen as a spread, a two-sided document puts two gardens side by side, each on the outside
of its page. They do not meet: the inner margins stay empty, and each garden has its own
seed, its own root and its own header.
\end{remark}

\subsection{Checking that nothing else moved}

A change like this one should leave every existing document exactly as it was, and that
can be checked without looking. The drawing code runs outside \TeX\ as well, in a small
Lua harness that replaces the three or four \TeX\ functions it calls, and writes the
content stream of a page to a file. Forty-five pages, five seeds with and without the space
under the text, in summer, in spring, in autumn with marks and without visitors, plus the
authored composition, were written before the change and again after it, and compared
byte for byte. They are identical: a page whose wider margin is on the left goes through
the same arithmetic in the same order as before, and the identity transform it uses is
the one it always used.

The growth animation is checked the same way. It replays every stem and every motif at
smaller sizes before the real one, restoring the random streams in between, and its last
frame has to be the page, byte for byte. It still is, for mirrored pages and for the theme
on the next page alike.

\newpage
\florilegiumtheme{geometric}
\section{Another garden}

From this page on, the garden grows from another set of primitives. The command
\code{\textbackslash florilegiumtheme\{geometric\}} leaves a mark in the text, like a
heading, so the theme changes from the page on which it stands; as an option of the
package it applies to the whole document:
\begin{center}\code{\textbackslash usepackage[grown,theme=geometric]\{florilegium\}}\end{center}

The growth code never knew what it was drawing. It reads a table of motifs, each with the
size of the disc it occupies, a weight and a function that draws it, and calls every
primitive by name. A theme is a Lua file that replaces some of those primitives, adds
rows to the table, or brings a composition of its own. Two come with the package:
\begin{itemize}
\item \emph{geometric}, on this page: the same plant cut from flat shapes. The stems run
straight between their nodes, the leaves are kites and toothed polygons folded into a
light and a dark facet, the flowers are rosettes of rhombi and star polygons, and
crystals grow among them.
\item \emph{math}, which is no plant at all: figures, symbols and equations cascading down
the margin and along the bottom, a large figure in the corner, and families of figures
that grow as they go, the Koch snowflake gaining a level, the cube truncating into the
octahedron.
\end{itemize}
Two options turn any margin down: \code{density} thins it, to a single vine at
$0.05$, and \code{intensity} pales it. The text ends here, and
\code{\textbackslash florilegiumtextend} tells the garden so: it grows into the space below
these lines as well.
\florilegiumtextend
\end{document}
